% !TEX root = ../tex/8-5.tex
% 第8章：本文は青、補足は緑
\newcommand{\asmain}{%
    \setbeamercolor{frametitle}{fg=white,bg=main}%
    \setbeamercolor{structure}{fg=main}%
    \setbeamercolor{itemize item}{fg=main}%
    \setbeamercolor{itemize subitem}{fg=main}%
    \setbeamercolor{enumerate item}{fg=main}%
}
\newcommand{\assuppl}{%
    \setbeamercolor{frametitle}{fg=white,bg=examplecol}%
    \setbeamercolor{structure}{fg=examplecol}%
    \setbeamercolor{itemize item}{fg=examplecol}%
    \setbeamercolor{itemize subitem}{fg=examplecol}%
    \setbeamercolor{enumerate item}{fg=examplecol}%
}
\newcommand{\supplmark}{%
    \setlength{\fboxsep}{1.5pt}%
    \setlength{\fboxrule}{0.7pt}%
    \fbox{\textbf{補足}}%
}
\asmain


%=============================================
\begin{frame}{8.5 From RL to DPO for Model Alignment}
    \begin{block}{この節の主題（論文 8.5）}
        人間の好みに合わせて生成モデルを揃える（alignment）。\\
        主流の RLHF は強力だが多段で不安定。\\
        \textbf{DPO} は同じ目標を、明示的報酬モデルと RL なしで狙う。
    \end{block}
    \begin{itemize}
        \item 言語モデル向け DPO（Rafailov et al., 2023）
        \item 拡散への拡張：\textbf{Diffusion-DPO}（Wallace et al., 2024）
        \item ゼミでは直観と鍵式を中心に追う（細部の導出は補足）
    \end{itemize}
\end{frame}
%=============================================
\asmain
\begin{frame}{8.5 の道筋}
    \begin{enumerate}
        \item \textbf{8.5.1 動機}\\
        RLHF の落とし穴と「直接好みで最適化できないか」
        \item \textbf{8.5.2 RLHF（Bradley--Terry 視点）}\\
        報酬モデルと KL 正則化付き方策最適化
        \item \textbf{8.5.3 DPO 枠組み}\\
        RLHF 目的の閉形式解 $\to$ 暗黙報酬 $\to$ ロジスティック損失
        \item \textbf{8.5.4 Diffusion-DPO}\\
        軌跡／逆過程への持ち上げと実用的代理損失
    \end{enumerate}
\end{frame}
%=============================================
\asmain
\begin{frame}{前回までとの接続}
    \begin{itemize}
        \item \textbf{8.1--8.3}：条件付きスコアの分解と CG / CFG
        \item \textbf{8.4}：学習なし誘導（タスク損失から誘導方向）
        \item \textbf{ここ（8.5）}：条件を満たしても好みから外れる問題へ\\
        \emph{control}（条件適合）と \emph{alignment}（人間好み）を区別
    \end{itemize}
    \begin{exampleblock}{好みによる傾き（プロローグのイメージ）}
        \[
            \tilde p_0(\vx_0\mid c)
            \propto
            p_0(\vx_0\mid c)\,
            \exp\!\bigl(\lambda_{\mathrm{pref}}\, r(\vx_0,c)\bigr)
        \]
        大きい $r$ ほど好ましいクリーン出力。KL 正則化では $\lambda_{\mathrm{pref}}=1/\beta$。
    \end{exampleblock}
\end{frame}
%=============================================
\section{8.5.1 Motivation}
\asmain
\begin{frame}{区切り：8.5.1 RLHF の落とし穴を避ける}
    \tableofcontents[currentsection]
\end{frame}
%=============================================
\asmain
\begin{frame}{Alignment の目標}
    ベース（例：SFT 済み）モデルを、\\
    人間が好む出力へ舵取りすること。
    \begin{block}{RLHF の三段パイプライン}
        \begin{enumerate}
            \item \textbf{SFT}：プロンプト--応答ペアで教師あり微調整
            \item \textbf{報酬モデリング（RM）}：好みペア $(c,\vx^w,\vx^l)$ から\\
            スカラー $r(c,\vx)$ を学習（勝者が敗者より高得点）
            \item \textbf{RL 微調整}：PPO 等で期待報酬を最大化しつつ、\\
            参照（SFT）分布への KL ペナルティで離れすぎを防ぐ
        \end{enumerate}
    \end{block}
\end{frame}
%=============================================
\asmain
\begin{frame}{RLHF が抱える実務上の欠点}
    \begin{itemize}
        \item \textbf{不安定・高コスト}：on-policy。各更新で現行方策から\\
        新しいサンプルを生成し直す必要がある
        \item \textbf{多モデル運用}：SFT・報酬・（しばしば）価値モデルを\\
        同時に学習・ホストする
        \item \textbf{代理目的}：最適化しているのは人間好みそのものではなく\\
        学習した報酬の代理。報酬の欠陥を突く（reward hacking）余地
    \end{itemize}
    \begin{alertblock}{Question 8.5.1}
        明示的報酬モデリングと不安定な RL 段階をやめ、\\
        好みデータ上でモデルを\textbf{直接}最適化できるか？
    \end{alertblock}
\end{frame}
%=============================================
\asmain
\begin{frame}{DPO の一言答え}
    \begin{block}{Direct Preference Optimization}
        多段 RLHF を、\textbf{単一の教師あり風ステップ}に置き換える。\\
        別報酬モデルも PPO も使わず、\\
        固定参照モデルに近づきつつ、好みペアへロジスティック損失で適合。
    \end{block}
    \begin{itemize}
        \item 鍵：KL 正則化付き RLHF 目的を書き換えると、\\
        方策と参照の\textbf{対数尤度比}が暗黙の報酬になる
        \item 同じ「参照への正則化」を保ちつつ、\\
        高価なロールアウトと明示的 RM を避けられる
    \end{itemize}
\end{frame}
%=============================================
\asmain
\begin{frame}{この節の残りでやること}
    \begin{enumerate}
        \item \textbf{8.5.2}：RLHF パイプラインと大規模報酬モデルへの依存を短く復習
        \item \textbf{8.5.3}：言語モデル向け DPO --- RM と RL を迂回する仕組み
        \item \textbf{8.5.4}：拡散への拡張（Diffusion-DPO）---\\
        生成モデリング設定での実用的・安定な alignment
    \end{enumerate}
    \vfill
    \footnotesize
    脚注の「方策」：プロンプト／履歴（状態）を応答／行動の分布へ写す写像。
\end{frame}
%=============================================
\section{8.5.2 RLHF: Bradley--Terry View}
\asmain
\begin{frame}{区切り：8.5.2 RLHF（Bradley--Terry 視点）}
    \tableofcontents[currentsection]
\end{frame}
%=============================================
\asmain
\begin{frame}{報酬モデル：好みの「審判」を学ぶ}
    同じプロンプト $c$ に対する候補応答へ、スカラー好み得点を付ける $r_{\bm{\psi}}$。
    \begin{itemize}
        \item データ $\mathcal{D}$：組 $(\tilde{\vx},\vx)$ とラベル $y$\\
        （$\tilde{\vx}$ が $\vx$ より好まれるか）
        \item $y\in\{0,1\}$（二値）または複数評価者の集約 $y\in[0,1]$（ソフト）
    \end{itemize}
    \begin{block}{ロジスティック損失（一般形）}
        \[
            \mathcal{L}_{\mathrm{RM}}(\bm{\psi})
            =
            -\mathbb{E}
            \Bigl[
            y\log\sigma\!\bigl(\Delta r\bigr)
            +(1-y)\log\bigl(1-\sigma(\Delta r)\bigr)
            \Bigr]
        \]
        ただし $\Delta r:=r_{\bm{\psi}}(c,\tilde{\vx})-r_{\bm{\psi}}(c,\vx)$，
        $\sigma(u)=1/(1+e^{-u})$。
    \end{block}
\end{frame}
%=============================================
\asmain
\begin{frame}{勝者・敗者の順序付け表記}
    実務では好みペアを $(\vx^w,\vx^l,c)$ の順序付きで保存することが多い。
    \begin{itemize}
        \item この慣例では常に $y=1$（勝者が先）
        \item $\tilde{\vx}=\vx^w$，$\vx=\vx^l$ と置けば一般形は次に退化：
    \end{itemize}
    \begin{block}{順序付き RM 損失（式 8.5.2）}
        \[
            \mathcal{L}_{\mathrm{RM}}(\bm{\psi})
            =
            -\mathbb{E}_{(\,c,\vx^w,\vx^l)\sim\mathcal{D}}
            \log\sigma\!\bigl(
            r_{\bm{\psi}}(c,\vx^w)-r_{\bm{\psi}}(c,\vx^l)
            \bigr)
        \]
    \end{block}
    \begin{exampleblock}{データの出どころ}
        人手キュレーション、異なるチェックポイント、\\
        事前学習済み条件付き拡散の生成、など。
    \end{exampleblock}
\end{frame}
%=============================================
\asmain
\begin{frame}{Bradley--Terry：スコア差 $\to$ 勝率}
    \[
        p_{r_{\bm{\psi}}}(\tilde{\vx}\succ\vx\mid c)
        :=
        \sigma\!\bigl(r_{\bm{\psi}}(c,\tilde{\vx})-r_{\bm{\psi}}(c,\vx)\bigr)
    \]
    \begin{itemize}
        \item \textbf{差だけが効く}：$r_{\bm{\psi}}(c,\cdot)$ は定数シフト不変
        \item 損失は「予測勝率」を人間ラベルに合わせ、\\
        勝者のスコアを敗者より高く押し上げる
    \end{itemize}
    \begin{alertblock}{直感（勾配）}
        $\Delta r\leftarrow\Delta r-\eta\bigl(\sigma(\Delta r)-y\bigr)$。\\
        $y=1$ なら $\Delta r$ が増える（勝者↑・敗者↓）。
    \end{alertblock}
\end{frame}
%=============================================
\assuppl
\begin{frame}{\supplmark：RM 損失 $=$ ベルヌーイ KL}
    1 サンプルのロジスティック損失は、観測ラベルと予測勝率の交差エントロピー：
    \[
        -\bigl[y\log p+(1-y)\log(1-p)\bigr]
        =
        D_{\mathrm{KL}}\!\bigl(\mathrm{Bern}(y)\,\|\,\mathrm{Bern}(p)\bigr)
        +H\!\bigl(\mathrm{Bern}(y)\bigr).
    \]
    \begin{block}{データ平均（式 8.5.3）}
        \[
            \mathcal{L}_{\mathrm{RM}}
            =
            \mathbb{E}_{\mathcal{D}}
            D_{\mathrm{KL}}\!\bigl(\mathrm{Bern}(y)\,\|\,\mathrm{Bern}(p_{r})\bigr)
            +
            \underbrace{\mathbb{E}_{\mathcal{D}} H(\mathrm{Bern}(y))}_{\bm{\psi}\text{に非依存}}
        \]
    \end{block}
    最小化 $=$ 人間の経験的勝率に予測勝率を合わせること（ソフトラベルでは定数オフセット差）。
\end{frame}
%=============================================
\asmain
\begin{frame}{固定報酬のもとでの KL 正則化方策最適化}
    フィット済み報酬 $r:=r_{\bm{\psi}^\star}$ と、参照 $\pi_{\mathrm{ref}}(\vx\mid c):=p_{\bm{\phi}^\star}(\vx\mid c)$（事前学習条件付き拡散）を固定。
    \begin{block}{RLHF 目的（式 8.5.4）}
        \[
            \max_{\bm{\theta}}
            \mathbb{E}_{c}
            \Bigl[
            \mathbb{E}_{\vx\sim\pi_{\bm{\theta}}(\cdot\mid c)}
            r(c,\vx)
            -
            \beta\,
            D_{\mathrm{KL}}\!\bigl(
            \pi_{\bm{\theta}}(\cdot\mid c)\,\|\,\pi_{\mathrm{ref}}(\cdot\mid c)
            \bigr)
            \Bigr]
        \]
    \end{block}
    \begin{itemize}
        \item 審判が高得点を付けるサンプルを増やす
        \item ただし参照（事前学習）から離れすぎない
    \end{itemize}
\end{frame}
%=============================================
\asmain
\begin{frame}{RLHF の二段まとめ}
    \begin{enumerate}
        \item 式 8.5.2（＝二値 KL の期待）を最小化して報酬 $r^\star$ をフィット
        \item 式 8.5.4 を解いて方策 $\pi^\star$ を最適化
    \end{enumerate}
    \begin{exampleblock}{補足}
        RM は\textbf{参照モデルが生成したペア}を必須としない。\\
        ただし、使う領域に近いモデルから集めると分布シフトが減り、\\
        報酬の信頼性が上がりやすい。
    \end{exampleblock}
    \begin{alertblock}{次へ}
        この KL 正則化目的には閉形式解があり、\\
        それを逆に読むと「明示的 $r$ なし」の DPO に到達する。
    \end{alertblock}
\end{frame}
%=============================================
\section{8.5.3 DPO Framework}
\asmain
\begin{frame}{区切り：8.5.3 DPO 枠組み}
    \tableofcontents[currentsection]
\end{frame}
%=============================================
\asmain
\begin{frame}{RLHF から DPO への橋：閉形式解}
    各プロンプト $c$ について、式 8.5.4 の最適方策はエネルギー型（式 8.5.5）：
    \[
        \pi^\star(\vx\mid c)
        =
        \frac{1}{Z(c)}\,
        \pi_{\mathrm{ref}}(\vx\mid c)\,
        \exp\!\bigl(r(c,\vx)/\beta\bigr).
    \]
    \begin{columns}[T,onlytextwidth]
        \begin{column}{0.48\textwidth}
            \begin{block}{小さい $\beta$}
                $\exp(r/\beta)$ が尖る\\
                $\to$ 高報酬に集中\\
                参照から遠く、多様性↓\\
                不安定・hack のリスク↑
            \end{block}
        \end{column}
        \begin{column}{0.48\textwidth}
            \begin{block}{大きい $\beta$}
                $\exp(r/\beta)$ が平坦\\
                $\to$ 参照に近い\\
                更新は保守的、多様性は参照寄り\\
                報酬の伸びは限定的
            \end{block}
        \end{column}
    \end{columns}
\end{frame}
%=============================================
\asmain
\begin{frame}{暗黙報酬：式 8.5.5 の反転}
    任意の方策 $\pi$（サポートが $\pi_{\mathrm{ref}}$ に含まれる）に対し
    \[
        r_{\pi}(c,\vx)
        =
        \beta\log\frac{\pi(\vx\mid c)}{\pi_{\mathrm{ref}}(\vx\mid c)}
        +\beta\log Z(c)
        \qquad\text{（式 8.5.6）}
    \]
    と置けば、その $\pi$ は「報酬 $r_{\pi}$ に対する」式 8.5.4 の最適解になる。
    \begin{itemize}
        \item $r_{\pi}$ は方策が誘起する\textbf{暗黙報酬}
        \item 定数 $\beta\log Z(c)$ はプロンプト依存だが、\\
        ペア比較（BT）では\textbf{打ち消される}
    \end{itemize}
\end{frame}
%=============================================
\asmain
\begin{frame}{ペアでは正規化定数が消える}
    \[
        r_{\pi}(c,\vx^w)-r_{\pi}(c,\vx^l)
        =
        \beta\log\frac{\pi(\vx^w\mid c)}{\pi_{\mathrm{ref}}(\vx^w\mid c)}
        -
        \beta\log\frac{\pi(\vx^l\mid c)}{\pi_{\mathrm{ref}}(\vx^l\mid c)}
    \]
    \begin{exampleblock}{要点}
        好み学習に必要なのは絶対スコアではなく、\\
        「参照に対する対数尤度の改善」の勝者--敗者差だけ。\\
        これが DPO 損失へ直結する。
    \end{exampleblock}
\end{frame}
%=============================================
\asmain
\begin{frame}{DPO の学習損失}
    暗黙報酬を BT（式 8.5.2）へ代入すると、$Z(c)$ が消え、
    \begin{block}{DPO 損失}
        \[
            \mathcal{L}_{\mathrm{DPO}}(\bm{\theta};\pi_{\mathrm{ref}})
            =
            -\mathbb{E}
            \log\sigma\!\Biggl(
            \beta\log\frac{\pi_{\bm{\theta}}(\vx^w\mid c)}{\pi_{\mathrm{ref}}(\vx^w\mid c)}
            -
            \beta\log\frac{\pi_{\bm{\theta}}(\vx^l\mid c)}{\pi_{\mathrm{ref}}(\vx^l\mid c)}
            \Biggr)
        \]
    \end{block}
    \begin{itemize}
        \item 勝者の（温度付き）優位を、参照からの対数比改善の差として押し上げる
        \item \textbf{単一の分類風微調整}で RLHF の目標に到達
        \item 明示的報酬モデルも、不安定な RL ロールアウトも不要
    \end{itemize}
\end{frame}
%=============================================
\asmain
\begin{frame}{DPO の読み方（一文）}
    \[
        -\log\sigma\!\Bigl(
        \beta\cdot
        \underbrace{
        \bigl[\text{log-ratio of }\pi_{\bm{\theta}}/\pi_{\mathrm{ref}}\bigr]_{\vx^w}
        -
        \bigl[\text{log-ratio}\bigr]_{\vx^l}
        }_{\text{winner--loser margin}}
        \Bigr)
    \]
    \begin{alertblock}{標語}
        「好みペア上のロジスティック回帰」として、\\
        方策の参照相対対数尤度を直接動かす。
    \end{alertblock}
\end{frame}
%=============================================
\assuppl
\begin{frame}{\supplmark：RLHF と DPO の対応表}
    \renewcommand{\arraystretch}{1.25}
    \footnotesize
    \begin{tabular}{lll}
        & \textbf{RLHF} & \textbf{DPO} \\
        \hline
        報酬 & 明示的 $r_{\bm{\psi}}$ を学習 & 方策の対数比が暗黙報酬 \\
        最適化 & PPO 等の RL（on-policy） & 単一の教師あり風損失 \\
        正則化 & $\beta$ KL を明示 & 同じ $\beta$ が温度として残る \\
        データ & 好みペア（＋ロールアウト） & 好みペアのみ \\
        狙い & 式 8.5.4 の解 & 同じ解クラスを直接フィット \\
    \end{tabular}
    \vfill
    \begin{exampleblock}{注意}
        DPO は「RL が不要」であって「好みデータが不要」ではない。\\
        参照モデルの選択と $\beta$ は、依然として重要な設計レバー。
    \end{exampleblock}
\end{frame}
%=============================================
\section{8.5.4 Diffusion-DPO}
\asmain
\begin{frame}{区切り：8.5.4 Diffusion-DPO}
    \tableofcontents[currentsection]
\end{frame}
%=============================================
\asmain
\begin{frame}{拡散では素朴な DPO がつらい}
    表記は 8.5.3 を引き継ぐ：$\pi_{\bm{\theta}}$ / $\pi_{\mathrm{ref}}$ は学習・参照の拡散方策\\
    （経路・遷移分布を含む）。前向きノイズ $p_t(\vx_t\mid\vx_0)$ は固定。
    \begin{alertblock}{なぜそのままではダメか}
        \begin{itemize}
            \item サンプル尤度 $\pi_{\bm{\theta}}(\vx\mid c)$ の評価には、\\
            ODE の瞬間的変数変換（ドリフトの発散）等が必要で重い
            \item サンプリング軌跡全体を通した微分は、\\
            勾配の消失・爆発のリスク
        \end{itemize}
    \end{alertblock}
    $\Rightarrow$ Diffusion-DPO は\textbf{経路レベル}で議論する\\
    （離散時間 DDPM を例に；連続時間も同様）。
\end{frame}
%=============================================
\asmain
\begin{frame}{経路と端点報酬}
    軌跡 $\vx_{0:T}:=(\vx_T,\ldots,\vx_0)$。\\
    逆時間マルコフ鎖の条件付き $\pi(\vx_{t-1}\mid\vx_t,c)$。\\
    $\vx_T$ は最高ノイズ、$\vx_0$ はクリーン出力。
    \begin{itemize}
        \item 人間の好みは最終出力上の $r(c,\vx_0)$ で定義される
        \item $\vx_0$ はどの軌跡にも含まれるので、\\
        「報酬の定義」自体を経路空間で変える必要はない
        \item 有用なのは、\textbf{逆過程全体}を正則化すること
    \end{itemize}
    \begin{exampleblock}{KL の連鎖則}
        経路分布の KL $\ge$ 端点周辺の KL。\\
        $\Rightarrow$ 潜在の denoising 軌跡の逸脱も制御する動機。
    \end{exampleblock}
\end{frame}
%=============================================
\asmain
\begin{frame}{経路版の KL 正則化目的}
    \[
        \max_{\bm{\theta}}
        \mathbb{E}_{c}
        \Bigl[
        \mathbb{E}_{\vx_{0:T}\sim\pi_{\bm{\theta}}(\cdot\mid c)}
        r(c,\vx_0)
        -
        \beta\,
        D_{\mathrm{KL}}\!\bigl(
        \pi_{\bm{\theta}}(\cdot\mid c)\,\|\,\pi_{\mathrm{ref}}(\cdot\mid c)
        \bigr)
        \Bigr]
    \]
    \begin{itemize}
        \item 報酬項は端点のまま
        \item KL は経路分布どうし（潜在軌跡のずれも含む）
        \item $\pi_{\bm{\theta}}(\cdot\mid c)$，$\pi_{\mathrm{ref}}(\cdot\mid c)$ は対応する\textbf{経路分布}
    \end{itemize}
\end{frame}
%=============================================
\asmain
\begin{frame}{経路の最適解と暗黙経路報酬}
    正規化定数が有限なら、各 $c$ で（式 8.5.7）
    \[
        \pi^\star(\vx_{0:T}\mid c)
        =
        \frac{1}{Z(c)}\,
        \pi_{\mathrm{ref}}(\vx_{0:T}\mid c)\,
        \exp\!\bigl(r(c,\vx_0)/\beta\bigr).
    \]
    反転して任意の $\pi$ に暗黙経路報酬を定義：
    \[
        R_{\pi}(c,\vx_{0:T})
        :=
        \beta\log\frac{\pi(\vx_{0:T}\mid c)}{\pi_{\mathrm{ref}}(\vx_{0:T}\mid c)}
        +\beta\log Z(c).
    \]
    \begin{itemize}
        \item 経路最適では $R_{\pi^\star}(c,\vx_{0:T})=r(c,\vx_0)$
        \item 一般の $\pi$ では潜在軌跡にも依存しうる
        \item 好みデータはクリーン端点だけを比較 $\Rightarrow$\\
        端点ごとに軌跡平均を取る
    \end{itemize}
\end{frame}
%=============================================
\asmain
\begin{frame}{端点ペアから Diff-DPO マージンへ}
    ラベル付き端点対 $(\vx_0^w,\vx_0^l)$ に対し、暗黙経路報酬の条件付き平均を比較：
    \[
        \mathcal{L}_{\mathrm{Diff\text{-}DPO}}
        =
        -\mathbb{E}
        \log\sigma\!\bigl(\Delta R(c;\bm{\theta})\bigr)
    \]
    \[
        \Delta R
        :=
        \mathbb{E}_{\vx_{1:T}^w\sim\pi_{\bm{\theta}}(\cdot\mid\vx_0^w,c)}
        \bigl[R_{\pi_{\bm{\theta}}}\bigr]
        -
        \mathbb{E}_{\vx_{1:T}^l\sim\pi_{\bm{\theta}}(\cdot\mid\vx_0^l,c)}
        \bigl[R_{\pi_{\bm{\theta}}}\bigr]
    \]
    \begin{exampleblock}{条件付き経路の意味}
        固定端点 $\vx_0$ のもと、その $\vx_0$ を生みうる逆時間軌跡の\\
        モデル誘導事後（拡散ブリッジ的）で平均する。
    \end{exampleblock}
\end{frame}
%=============================================
\asmain
\begin{frame}{式 8.5.8 が実務でつらい三理由}
    \begin{enumerate}
        \item \textbf{端点条件付き経路事後が intractable}\\
        サンプラは $\vx_T\to\cdots\to\vx_0$（無制約）。\\
        端点固定は一般に閉形式がなく、サンプリングも高価
        \item \textbf{入れ子で $\bm{\theta}$ 結合の期待}\\
        経路分布も被積分関数 $R$ も $\bm{\theta}$ 依存\\
        $\Rightarrow$ REINFORCE / pathwise 結合、高分散勾配
        \item \textbf{長い鎖と重い逆伝播}\\
        経路対数密度は $O(T)$ ステップ（$T\sim 10^2$--$10^3$）。\\
        winner/loser・policy/ref の両方で計算し、予算を超えやすい
    \end{enumerate}
\end{frame}
%=============================================
\asmain
\begin{frame}{Jensen で上界：経路尤度比へ}
    $-\log\sigma(\cdot)$ は凸。内側期待を外へ出すと上界：
    \[
        \mathcal{L}_{\mathrm{Diff\text{-}DPO}}
        \le
        -\mathbb{E}
        \log\sigma\!\Bigl(
        R_{\pi_{\bm{\theta}}}(c,\vx_{0:T}^w)
        -
        R_{\pi_{\bm{\theta}}}(c,\vx_{0:T}^l)
        \Bigr)
    \]
    暗黙報酬の恒等式と定数打ち消しのあと（式 8.5.9）：
    \[
        \le
        -\mathbb{E}
        \log\sigma\!\Biggl(
        \beta\log\frac{\pi_{\bm{\theta}}(\vx_{0:T}^w\mid c)}{\pi_{\mathrm{ref}}(\vx_{0:T}^w\mid c)}
        -
        \beta\log\frac{\pi_{\bm{\theta}}(\vx_{0:T}^l\mid c)}{\pi_{\mathrm{ref}}(\vx_{0:T}^l\mid c)}
        \Biggr).
    \]
    \begin{alertblock}{まだ}
        経路全体の尤度比が残る $\Rightarrow$ マルコフ分解で一歩ずつへ。
    \end{alertblock}
\end{frame}
%=============================================
\asmain
\begin{frame}{マルコフ分解：経路比 $\to$ ステップ和}
    \[
        \pi_{\bm{\theta}}(\vx_{0:T}\mid c)
        =
        \pi_{\bm{\theta}}(\vx_T\mid c)
        \prod_{t=1}^{T}
        \pi_{\bm{\theta}}(\vx_{t-1}\mid\vx_t,c)
    \]
    参照も同様。時刻 $T$ の事前が共通なら先頭比は 1 で、
    \[
        \log\frac{\pi_{\bm{\theta}}(\vx_{0:T}\mid c)}{\pi_{\mathrm{ref}}(\vx_{0:T}\mid c)}
        =
        \sum_{t=1}^{T}
        \log\frac{\pi_{\bm{\theta}}(\vx_{t-1}\mid\vx_t,c)}{\pi_{\mathrm{ref}}(\vx_{t-1}\mid\vx_t,c)}.
    \]
    したがって上界は
    \[
        -\mathbb{E}
        \log\sigma\!\Bigl(
        \beta\sum_{t=1}^{T}\Delta_t
        \Bigr)
    \]
    の形（各 $\Delta_t$ は winner/loser の 1 ステップ対数比差）。
\end{frame}
%=============================================
\asmain
\begin{frame}{時刻の一様サンプルとさらなる Jensen}
    $\beta\sum_s\Delta_s=\mathbb{E}_{t\sim\mathrm{Unif}\{1,\ldots,T\}}[\beta T\,\Delta_t]$。\\
    再び $-\log\sigma$ の凸性より
    \[
        \mathcal{L}_{\mathrm{Diff\text{-}DPO}}
        \le
        -\mathbb{E}
        \log\sigma\!\bigl(\beta T\,\Delta_t\bigr).
    \]
    \begin{itemize}
        \item 期待は $(c,\vx_0^w,\vx_0^l)$、$t$、および端点条件付き遷移対の上
        \item \textbf{1 遷移}だけに依存するが、\\
        端点条件付き遷移のサンプリングは依然 intractable
    \end{itemize}
\end{frame}
%=============================================
\asmain
\begin{frame}{前向きノイズによる実用近似}
    Diffusion-DPO では、モデル誘導の条件付き経路を\\
    \textbf{固定の前向きノイズ過程}で近似する。
    \begin{enumerate}
        \item クリーン端点 $\vx_0$ から $\vx_t\sim p_t(\vx_t\mid\vx_0)$ をサンプル
        \item $(\vx_0,\vx_t)$ 条件下の前向き事後 $p(\vx_{t-1}\mid\vx_t,\vx_0)$ は tractable
        \item その事後で逆遷移の対数比を平均
    \end{enumerate}
    \[
        \mathbb{E}_{p(\vx_{t-1}\mid\vx_t,\vx_0)}
        \log\frac{\pi_{\bm{\theta}}(\vx_{t-1}\mid\vx_t,c)}{\pi_{\mathrm{ref}}(\vx_{t-1}\mid\vx_t,c)}
        =
        -D_{\mathrm{KL}}(p\|\pi_{\bm{\theta}})
        +D_{\mathrm{KL}}(p\|\pi_{\mathrm{ref}}).
    \]
\end{frame}
%=============================================
\asmain
\begin{frame}{ガウス逆核 $\to$ DSM 風 MSE 差}
    逆分散が共通のガウス核では、KL は重み付き denoising 回帰誤差へ退化。\\
    $\epsilon$-予測では（式 8.5.10）
    \[
        \mathbb{E}
        \log\frac{\pi_{\bm{\theta}}}{\pi_{\mathrm{ref}}}
        =
        -\lambda_t
        \Bigl(
        \|\hat{\veps}_{\bm{\theta}}(\vx_t,t,c)-\veps\|^2
        -
        \|\hat{\veps}_{\mathrm{ref}}(\vx_t,t,c)-\veps\|^2
        \Bigr),
    \]
    ただし $\vx_t=\alpha_t\vx_0+\sigma_t\veps$，$\veps\sim\mathcal{N}(\bm{0},I)$，$\lambda_t>0$。
    \begin{block}{略記}
        \[
            \Delta\mathrm{MSE}(\vx_t;\veps)
            :=
            \|\hat{\veps}_{\bm{\theta}}-\veps\|^2
            -
            \|\hat{\veps}_{\mathrm{ref}}-\veps\|^2
        \]
    \end{block}
\end{frame}
%=============================================
\asmain
\begin{frame}{実用的な 1 ステップ Diffusion-DPO 代理損失}
    \[
        \widetilde{\mathcal{L}}_{\mathrm{Diff\text{-}DPO}}
        =
        -\mathbb{E}
        \log\sigma\!\Bigl(
        \beta T\lambda_t
        \bigl(
        \Delta\mathrm{MSE}(\vx_t^l;\veps^l)
        -
        \Delta\mathrm{MSE}(\vx_t^w;\veps^w)
        \bigr)
        \Bigr)
    \]
    期待は $(c,\vx_0^w,\vx_0^l)\sim\mathcal{D}$，$t\sim\mathrm{Unif}$，$\veps^w,\veps^l\sim\mathcal{N}(\bm{0},I)$。\\
    $\vx_t^{w/l}=\alpha_t\vx_0^{w/l}+\sigma_t\veps^{w/l}$。
    \begin{itemize}
        \item 共有ノイズ $\veps^w=\veps^l$ は分散低減の変種（形は DPO 風のまま）
        \item 参照は $\log\sigma$ 内のマージンに残り、例ごとの重み付けを形作る\\
        （参照自体への直接勾配はない）
    \end{itemize}
\end{frame}
%=============================================
\asmain
\begin{frame}{直感：何を押し上げているか}
    $\widetilde{\mathcal{L}}$ を最小化すると、参照正規化の\\
    \textbf{勝者--敗者マージン}が増える。
    \begin{block}{等価な読み}
        凍結参照に対し、\\
        勝者側の denoising 誤差をより減らし、\\
        敗者側より相対的に良く当てるよう促す。
    \end{block}
    \begin{exampleblock}{ゼミ向け要約}
        拡散では「サンプル尤度の DPO」を避け、\\
        経路 $\to$ ステップ $\to$ 前向き近似 $\to$ $\epsilon$-MSE 差、\\
        と tractable な代理へ落とす。
    \end{exampleblock}
\end{frame}
%=============================================
\assuppl
\begin{frame}{\supplmark：Diffusion-DPO の設計フロー（地図）}
    \footnotesize
    \begin{enumerate}
        \item 端点報酬＋経路 KL 正則化（RLHF の経路版）
        \item 閉形式 $\Rightarrow$ 暗黙経路報酬 $R_{\pi}$
        \item 端点条件付き平均で BT マージン（式 8.5.8）--- しかし intractable
        \item Jensen：経路尤度比の上界（式 8.5.9）
        \item マルコフ分解＋時刻一様化：1 ステップ対数比
        \item 前向き事後で近似：KL $=$ 重み付き MSE 差
        \item 実用損失：$\log\sigma$ 内の $\Delta\mathrm{MSE}$ 勝者--敗者差
    \end{enumerate}
\end{frame}
%=============================================
\asmain
\begin{frame}{8.5 のまとめ}
    \begin{itemize}
        \item RLHF は強いが、RM＋on-policy RL で重く不安定になりうる
        \item DPO：KL 正則化 RLHF の閉形式を反転し、\\
        対数比を暗黙報酬として好みペアへ直接フィット
        \item 拡散では尤度評価が重い $\Rightarrow$ 経路版へ持ち上げ、\\
        Jensen・マルコフ・前向き近似で DSM 風の代理損失へ
        \item 結果：明示的報酬モデルなしの、安定な alignment 微調整
    \end{itemize}
    \begin{alertblock}{次節}
        第8章全体の道具立てをまとめ、第9章（高速ソルバ）へ橋渡しする。
    \end{alertblock}
\end{frame}
